% thesis_main.tex - Master Thesis Template (IBF, Leibniz Universität Hannover)
\documentclass[12pt,a4paper,oneside]{report}

% --- Geometry (IBF requirements) ---
\usepackage[top=3cm, bottom=3cm, left=2cm, right=5cm]{geometry}

% --- Font: Times New Roman ---
\usepackage{mathptmx}
\usepackage[T1]{fontenc}
\usepackage[utf8]{inputenc}

% --- Spacing ---
\usepackage{setspace}
\onehalfspacing

% --- Language ---
\usepackage[english]{babel}

% --- Bibliography ---
\usepackage[round]{natbib}
\bibliographystyle{plainnat}

% --- Tables and Figures ---
\usepackage{booktabs}
\usepackage{graphicx}
\usepackage{caption}
\captionsetup{font={small}}

% --- Headings ---
\usepackage{titlesec}
\titleformat{\chapter}[display]
  {\normalfont\bfseries\fontsize{14}{16}\selectfont}
  {\chaptertitlename\ \thechapter}{12pt}
  {\fontsize{14}{16}\selectfont}
\titlespacing*{\chapter}{0pt}{0pt}{6pt}
\titleformat{\section}{\normalfont\bfseries\fontsize{12}{14}\selectfont}{\thesection}{1em}{}
\titlespacing*{\section}{0pt}{12pt}{6pt}
\titleformat{\subsection}{\normalfont\bfseries\fontsize{12}{14}\selectfont}{\thesubsection}{1em}{}
\titlespacing*{\subsection}{0pt}{12pt}{6pt}

% --- Page numbering bottom right ---
\usepackage{fancyhdr}
\pagestyle{fancy}
\fancyhf{}
\fancyfoot[R]{\thepage}
\renewcommand{\headrulewidth}{0pt}

% --- Hyperlinks ---
\usepackage{hyperref}
\hypersetup{colorlinks=false}

\begin{document}

% ============================================================
% FRONT MATTER (roman numerals)
% ============================================================
\pagenumbering{roman}

% --- Cover Sheet Placeholder ---
\begin{titlepage}
\centering
\vspace*{2cm}
{\Large Leibniz Universit\"at Hannover\\Institute of Banking and Finance\par}
\vspace{2cm}
{\LARGE\bfseries Textual Analysis of SEC 10-K Filings:\\
Replicating and Extending Loughran-McDonald (2011)\\
with Dictionary and FinBERT-Based Sentiment Measures\par}
\vspace{2cm}
{\large Master Thesis\par}
\vspace{1cm}
{\large Author Name\\Matriculation Number\par}
\vspace{1cm}
{\large Supervisor: Prof.\ Name\\Second Examiner: Prof.\ Name\par}
\vspace{1cm}
{\large Hannover, \today\par}
\end{titlepage}

% --- Table of Contents ---
\tableofcontents
\newpage
\listoffigures
\newpage
\listoftables
\newpage

% --- Abbreviations ---
\chapter*{List of Abbreviations}
\addcontentsline{toc}{chapter}{List of Abbreviations}
\begin{tabular}{ll}
CIK    & Central Index Key \\
CRSP   & Center for Research in Security Prices \\
EDGAR  & Electronic Data Gathering, Analysis, and Retrieval \\
FF4    & Fama-French Four-Factor Model \\
H4N    & Harvard-IV-4 Negative Word List \\
LM     & Loughran-McDonald \\
MD\&A  & Management's Discussion and Analysis \\
SEC    & Securities and Exchange Commission \\
SUE    & Standardized Unexpected Earnings \\
TF-IDF & Term Frequency-Inverse Document Frequency \\
\end{tabular}
\newpage

% ============================================================
% MAIN TEXT (arabic numerals)
% ============================================================
\pagenumbering{arabic}

% ────────────────────────────────────────────────────────────
\chapter{Introduction}
\label{ch:introduction}

\section{Motivation and Research Question}
\label{sec:motivation}
% FROM CODE: Replication spec (lm2011_replication_spec.yaml) targets LM2011 paper
%   "When Is a Liability Not a Liability?" with filing_sample_window [1994-01-01, 2008-12-31].
% FROM CODE: Extension spec (lm2011_extension_spec.yaml) extends sample end to 2024-12-31,
%   adds FinBERT comparison axis, and adds candidate item scopes (Item 1A, 7A).
% FROM CODE: Replication spec excludes Rule 10b-5 fraud labels and material-weakness labels.
% TODO: Motivate the research question from the literature (not from code).

\section{Contribution and Scope}
\label{sec:contribution}
% THREE RESEARCH AXES (from extension spec):
% (E1) Horizon Extension: Replicate LM2011 on 1994-2008, then extend identically to 2009-2024.
%      Report results separately for overlap and extension windows.
% (E2) Additional Item Scopes: Beyond full 10-K and MD&A (Item 7), evaluate Item 1A (Risk Factors,
%      available from Dec 2005) and potentially Item 7A (Market Risk, from Apr 1997).
% (E3) Dictionary vs. FinBERT Comparison: Run same Fama-MacBeth regressions with FinBERT
%      sentence-level sentiment aggregated to document level, alongside LM dictionary features.
% SCOPE BOUNDARIES: Excludes Rule 10b-5 fraud label regressions and material-weakness labels
%   from LM2011 (data not available). Excludes fine-tuning FinBERT (frozen model only).

\section{Thesis Structure}
\label{sec:structure}
% Ch.2 Literature Review: Dictionary vs. ML approaches, LM2011 findings, FinBERT architecture.
% Ch.3 Data and Sample: EDGAR universe, CRSP/Compustat linking, Refinitiv enrichment, descriptives.
% Ch.4 Item Extraction: Pre-processing, boundary detection, regime-aware validation, performance.
% Ch.5 Text Features: Dictionary proportional/TF-IDF features, FinBERT scoring, aggregation.
% Ch.6 Methodology: Excess returns, SUE, Fama-MacBeth regressions, trading strategy.
% Ch.7 Results: Replication (1994-2008), extension (2009-2024), FinBERT comparison.
% Ch.8 Discussion: Replication fidelity, temporal stability, dictionary vs. model trade-offs.
% Ch.9 Conclusion: Summary, implications, future research.


% ────────────────────────────────────────────────────────────
\chapter{Literature Review}
\label{ch:literature}

\section{Textual Analysis in Finance}
\label{sec:textual_analysis}

\subsection{Dictionary-Based Approaches}
\label{sec:dictionary_approaches}
% FROM CODE: Pipeline loads these dictionary files (lm2011_sample_post_refinitiv_runner.py):
%   LM categories: Fin-Neg.txt, Fin-Pos.txt, Fin-Unc.txt, Fin-Lit.txt, MW-Strong.txt, MW-Weak.txt.
%   Harvard: Harvard_IV_NEG_Inf.txt (H4N informal negative -- used as comparison baseline).
%   Trading strategy uses: fin_neg (financial negative subset) + h4n_inf.
% FROM CODE: Six LM categories in output: negative, positive, uncertainty, litigious, modal_strong, modal_weak.
% TODO: Describe H4N general-purpose vs. LM finance-specific distinction from the literature.
% TODO: Cite Tetlock (2007), Henry (2008), Li (2010), Loughran & McDonald (2011, 2016).

\subsection{Machine Learning and Transformer-Based Approaches}
\label{sec:ml_approaches}
% FROM CODE: Extension spec axis E3 positions FinBERT as "compact BERT-based finance model"
%   to compare against dictionary baseline. Model family: "BERT-based finance encoder".
% TODO: Trace evolution from bag-of-words -> TF-IDF -> embeddings -> transformers -> BERT.
% TODO: Cite Mikolov et al. (2013), Pennington et al. (2014), Vaswani et al. (2017), Devlin et al. (2019).

\section{Sentiment and Market Outcomes}
\label{sec:sentiment_markets}

\subsection{Loughran-McDonald (2011): Methodology and Findings}
\label{sec:lm2011_review}
% FROM CODE (replication spec + lm2011_regressions.py): Pipeline replicates LM2011 with:
%   - filing_sample_window: [1994-01-01, 2008-12-31].
%   - Included forms: 10-K, 10-K405 (raw form types). Amendments excluded.
%   - Stock filters: SHRCD in {10, 11}, EXCHCD in {1, 2, 3} (lm2011.py constants).
%   - Quarterly Fama-MacBeth regressions (run_lm2011_quarterly_fama_macbeth).
%   - Dependent variable: filing_period_excess_return (4-day event window [0,3]).
%   - Controls: log_size, log_book_to_market, pre_ffalpha, log_share_turnover,
%     nasdaq_dummy, institutional_ownership + FF48 industry dummies.
%   - Newey-West standard errors with nw_lags=1 (default).
%   - Tables implemented: IV (full 10-K), V (MD&A), VI (all 6 LM categories), VIII (SUE).
%   - Trading strategy: July 1997 - June 2007, quintile portfolios, 12-month holding.
% TODO: Summarize LM2011 key findings from the paper itself.

\subsection{Subsequent Empirical Extensions}
\label{sec:subsequent_extensions}
% TODO: Review and cite post-2011 studies using LM dictionaries or extending the sample.
% TODO: Cite Loughran & McDonald (2016), Bodnaruk et al. (2015), Garcia (2013), etc.

\section{FinBERT and Pre-Trained Language Models in Finance}
\label{sec:finbert_lit}

\subsection{Architecture and Training}
\label{sec:finbert_architecture}
% FROM CODE (extension spec): default_candidate = "FinBERT", model_family = "BERT-based finance encoder".
%   Extension spec design questions (still pending):
%   - Choose scoring grain: sentence, paragraph, item, or filing chunk.
%   - Define chunk aggregation from model outputs to doc_id-level and item-level features.
%   - Pin tokenizer, truncation, batching, and pooling rules for reproducibility.
%   - Decide whether to keep model frozen or allow fine-tuning stage.
% TODO: Describe FinBERT architecture, training data, and output format from the literature.
% TODO: Clarify which FinBERT variant (Araci 2019 vs. Huang et al. 2020).

\subsection{Sentence-Level vs.\ Document-Level Sentiment}
\label{sec:sentence_vs_doc}
% FinBERT processes single sentences (max 512 WordPiece tokens).
% A 10-K filing contains thousands of sentences -- aggregation required.
% Design choices for this thesis (from extension spec, status: planned):
%   - Scoring grain: sentence-level (one softmax per sentence).
%   - Aggregation to doc_id level: mean pooling of sentence scores as baseline.
%   - Alternative: token-count-weighted mean, or proportion of sentences above threshold.
%   - Truncation: sentences exceeding 512 tokens truncated (rare for well-split text).
%   - Must pin tokenizer, truncation, batching, and pooling rules for reproducibility.

\subsection{Comparison with Dictionary Methods}
\label{sec:dict_vs_model_lit}
% FROM CODE (extension spec): comparison_contract requires:
%   - Model-based features emitted alongside dictionary features on same sample and scope.
%   - Evaluation emphasizes incremental explanatory power and robustness, not only raw correlation.
%   - Both feature sets share doc_id, filing_date, text_scope grain.
% TODO: Review existing literature comparing dictionary and model-based text features.


% ────────────────────────────────────────────────────────────
\chapter{Data and Sample Construction}
\label{ch:data}

\section{SEC EDGAR Filing Universe}
\label{sec:edgar_universe}
% SOURCE: SEC EDGAR full-text archive, processed as yearly parquet files.
% INCLUDED FORMS: 10-K and 10-K405 (raw form types, not normalized).
%   10-K405 was the annual report form for accelerated filers until SEC eliminated it (~2003).
% EXCLUDED FORMS: All amendments (10-K/A, 10-K405/A), transition reports (10-KT, 10-KT/A),
%   small business forms (10-KSB, 10-KSB/A, 10KSB40, 10KSB40/A).
%   Rationale: amendments restate prior filings and would double-count text signals.
% SAMPLE WINDOWS:
%   - Replication: 1994-01-01 to 2008-12-31 (matching LM2011).
%   - Extension: 2009-01-01 to 2024-12-31.
%   - Full: 1994-2024.
% DEDUPLICATION: One filing per CIK per fiscal year. If multiple filings exist for the same
%   CIK-year pair, the earliest filing_date is retained. Minimum 180-day spacing between
%   consecutive filings from the same CIK to avoid fiscal-year-change duplicates.
% IDENTIFIER: doc_id = "{cik_10}:{accession}" where CIK is zero-padded to 10 digits.

\section{CRSP/Compustat Linking and Sample Filters}
\label{sec:crsp_ccm}
% CANONICAL LINK TABLE: Merges two CCM sources:
%   (1) linkhistory: GVKEY-PERMNO windows with LINKTYPE/LINKPRIM metadata (source_priority=1).
%   (2) linkfiscalperiodall: Broader coverage including sparse fallback rows (source_priority=2).
%   Combined with company history for GVKEY-CIK windows.
%   Valid date window = max(link_start, cik_start) to min(link_end, cik_end).
% LINK QUALITY TIERS (row_quality_tier):
%   Tier 10 (highest): PERMNO>0, valid window, not sparse, full link metadata.
%   Tier 20: PERMNO>0, valid window, not sparse, partial metadata.
%   Tier 40: PERMNO>0, no window, not sparse.
%   Tier 90 (lowest): Everything else.
% LINK QUALITY WEIGHTS: 4.0 (LINKTYPE LC/LU + LINKPRIM P), 3.0 (LINKPRIM P only),
%   2.0 (LINKTYPE LC/LU only), 0.1 (sparse fallback), 1.0 (default).
% STOCK FILTERS:
%   - SHRCD in {10, 11}: US ordinary common shares only.
%   - EXCHCD in {1, 2, 3}: NYSE, NYSE MKT (AMEX), NASDAQ only.
%   - Pre-filing stock price >= \$3.00 (penny stock exclusion).
% KEY IDENTIFIERS:
%   - GVKEY (firm-level, cast to Int32 before joins -- raw is string).
%   - PERMNO (security-level, one per share class).
%   - CIK (SEC filer identifier, 10-digit zero-padded).
% NOTE: A dual-class firm has one GVKEY but multiple PERMNOs. Joining on the wrong
%   key silently duplicates rows -- the canonical link table handles this.

\section{Refinitiv/LSEG Enrichment}
\label{sec:refinitiv}
% PURPOSE: Provides analyst forecast data (for SUE) and institutional ownership
%   that are not available in CRSP/Compustat.
% BRIDGE PIPELINE (4 stages):
%   (1) Bridge: Build universe mapping PERMNO+GVKEY to CUSIP/ISIN/TICKER identifiers.
%   (2) Lookup: Query LSEG API to resolve identifiers to RICs (Refinitiv Instrument Codes).
%   (3) Resolution: Determine effective_collection_ric per bridge row (priority: accepted_ric
%       from manual review > extended_ric borrowed from adjacent PERMNO/LIID rows > null).
%   (4) Instrument Authority: Final linkage table PERMNO+GVKEY -> effective_collection_ric.
% ANALYST DATA (via TR.EPSMean, TR.EPSActValue, TR.EPSStdDev, TR.EPSNumberOfEstimates):
%   - actual_eps: Reported EPS from earnings announcement.
%   - forecast_consensus_mean: Mean analyst EPS forecast.
%   - forecast_dispersion: Std. dev. of analyst estimates.
%   - forecast_revision_4m: Change in consensus over 4 months before announcement.
%   Used for SUE = (actual_eps - forecast_consensus_mean) / pre_filing_price.
% INSTITUTIONAL OWNERSHIP (via TR.CategoryOwnershipPct):
%   - Target quarter: most recent quarter-end strictly before filing_date.
%   - Category: "Holdings by Institutions" only.
%   - Fallback: If exact quarter missing, search window up to 180 days after quarter-end.
%   - Values clamped to [0, 100] range.
% LEDGER SYSTEM: SQLite-backed idempotent request tracking. States: PENDING, RUNNING,
%   SUCCEEDED, RETRYABLE_ERROR, DEFERRED_DAILY_LIMIT, FATAL_ERROR.
%   Stale-running detection (>900s) auto-requeues. Safe to re-run any stage.

\section{Sample Descriptive Statistics}
\label{sec:descriptives}
% TODO: Generate from actual data -- filing counts, sample attrition, summary statistics.
% FROM CODE: Key variables available in regression panels (lm2011_regressions.py):
%   filing_period_excess_return, sue, log_size, log_book_to_market, pre_ffalpha,
%   log_share_turnover, nasdaq_dummy, institutional_ownership, analyst_dispersion,
%   analyst_revisions, ff48_industry_id, token_count_full_10k, token_count_mda.
% FROM CODE: Sample filters that drive attrition (lm2011.py, transforms.py, lm2011_pipeline.py):
%   form type inclusion/exclusion -> CRSP/CCM join -> SHRCD/EXCHCD filter ->
%   pre_filing_price >= $3 -> min 60 days pre-filing data -> accounting data availability.


% ────────────────────────────────────────────────────────────
\chapter{Item Extraction from SEC 10-K Filings}
\label{ch:item_extraction}

\section{Pre-Processing Pipeline}
\label{sec:preprocessing}
% Four sequential transformations applied to each raw filing:
% (1) NEWLINE NORMALIZATION: Replace \r\n and \r with \n (mixed line endings from
%     different submission platforms across 30 years of EDGAR).
% (2) EDGAR HEADER REMOVAL: Strip <SEC-Header>, <Header>, <FileStats>, <XML_Chars> blocks
%     using tag-aware regex. Preserve filing_date, period-of-report, CIK, accession as
%     structured fields. Cross-validate header-derived identifiers against filename-encoded ones.
% (3) WRAPPED HEADING REPAIR: Rejoin headings split across lines by PDF/HTML conversion.
%     Patterns: "TABLE\nOF\nCONTENTS" -> "TABLE OF CONTENTS",
%     "PART\nII" -> "PART II", "ITEM\n1A" -> "ITEM 1A".
%     Conservative: only high-precision patterns corrected to avoid corrupting body text.
% (4) SPARSE LAYOUT DETECTION: Classify filing layout based on first 800 lines.
%     Sparse if ANY of: (a) max line length >= 8,000 chars;
%     (b) <= 200 lines AND max line length >= 2,000 chars;
%     (c) <= 300 lines AND avg line length >= 1,000 AND max >= 4,000 chars.
%     Sparse filings get mid-line heading scanning in boundary detection.
% Additionally: TOC regions detected via "Table of Contents" / "Index" markers, dot leaders
%   (sequences of 4+ periods followed by page numbers), and formatting patterns.
%   Lines in TOC ranges are masked (excluded from item boundary search).
%   Requires >= 4 valid ITEM entries with page numbers (1-500) to confirm TOC cutoff.

\section{Item Boundary Detection}
\label{sec:boundary_detection}
% Three-stage process: part markers -> item headings -> boundary selection.
% PART MARKERS: Line-start-anchored regex for "PART I" through "PART IV".
%   Classified as high-confidence (standalone heading) or low-confidence (embedded in text).
% ITEM HEADING DETECTION -- three complementary regex strategies:
%   (1) Line-start pattern (strictest): "ITEM" at/near line start, optionally with part prefix.
%       Primary matcher for standard-layout filings.
%   (2) Candidate pattern (permissive): "ITEM" at any position within a line.
%       Activated only for sparse-layout filings.
%   (3) Combined part-item pattern: "PART II -- ITEM 5. Market for..." on single line.
% CONFIDENCE LEVELS (integer 0-2):
%   HIGH (2): Item at line start + title matches SEC specification + not followed by prose.
%   MEDIUM (1): Midline item with title match, or standalone numeric dot heading.
%   LOW (0): Midline without title match, compound lines, or page-like context.
% CROSS-REFERENCE FILTERING: Lines starting with "See", "Refer to", "As discussed in",
%   "Pursuant to" rejected. Lines with dot leaders flagged as TOC entries.
%   Multi-item compound references ("Items 1, 1A, and 2") downgraded to low confidence.
% TWO-PASS STRATEGY:
%   Pass 1: Normal TOC masking + detected sparse layout setting.
%   Pass 2 (if pass 1 yields 0 items): Disable TOC masking, force sparse scanning on.
%   Ensures even atypical filings produce extractable items at cost of lower precision.
% BOUNDARY SELECTION: When multiple candidates for same item, prefer highest confidence,
%   break ties by position (later preferred -- earlier more likely TOC entries).

\section{Regime-Aware Item Validation}
\label{sec:regime}
% Form 10-K structure changed 12 times between 1997 and 2023 (SEC releases).
% KEY CHANGES (with effective dates from regime JSON):
%   1997 (33-7386): Item 7A Market Risk added (effective_date 1997-04-11).
%   2005 (33-8591): Items 1A Risk Factors and 1B Unresolved Comments added (2005-12-01).
%   2009 (33-9089): Item 4 "Voting Results" removed (2010-02-28).
%   2011 (33-9286): Item 4 repurposed as "Mine Safety Disclosures" (2012-01-27).
%   2020 (33-10890): Item 6 "Selected Financial Data" reserved (2021-02-10).
%   2023 (33-11216): Item 1C "Cybersecurity" added (fiscal years ending >= 2023-12-15).
% REGIME ENCODING: JSON document mapping each item to temporal validity windows.
%   Each window: start_date, end_date (or null=open), trigger type (effective_date or
%   fiscal_year_end_ge), and canonical label (e.g., "II:7_MDA").
%   Item 4 has THREE successive regimes: Voting -> Reserved -> Mine Safety.
% POST-EXTRACTION ANNOTATION: Three fields added per extracted item:
%   (1) canonical_item: standardized "Part:ItemNumber_Description" label.
%   (2) exists_by_regime: boolean -- is this item expected for this filing date/form?
%   (3) item_status: "active", "reserved", "optional", or "time_varying".
% Items not existing by regime can be dropped, preventing phantom items in text features.
% Critical for extension to 2024: Item 1C must be distinguished from items present since 1994.

\section{Quality Assurance and Extraction Performance}
\label{sec:extraction_quality}
% CONFIDENCE-BASED FILTERING: Each item carries confidence (0-2). Downstream consumers
%   can filter to trade coverage vs. precision.
% RESERVED-ITEM TRUNCATION: Items containing "[Reserved]" or "Reserved" detected via
%   pattern matching (also matches "None", "N/A", "Not applicable", "Not required").
%   Content truncated at reservation statement to prevent boilerplate inflating word counts.
% EMPTY-ITEM DETECTION: After extraction, items tested against whitespace/page-marker pattern.
%   Items of only whitespace, page numbers, or pagination artifacts flagged as empty.
% EMBEDDED HEADING DETECTION: Post-extraction scan for cross-references within item text
%   ("See Part II, Item 7"). Logged for review but do not trigger re-extraction.
% FAIL-FAST: Any extraction exception halts pipeline with full context (year, doc_id, CIK,
%   accession, form type, text length).
% CYTHON FAST PATH: Two core scanning functions (scan_part_markers, scan_item_boundaries)
%   reimplemented as C-extension via Cython. Uses typed memoryviews and inlined helpers.
%   3-10x faster than pure Python. Transparent fallback: if compiled extension unavailable,
%   Python version used automatically with one-time warning. Metrics track fast-path
%   successes, failures, and fallbacks for reproducibility auditing.
% FUZZ TESTING: Hypothesis framework generates 500 randomized document scenarios per suite.
%   Custom strategy produces realistic SEC document fragments (PART/ITEM headings, prose,
%   TOC entries, pagination). Parity assertion: Python and Cython produce identical results.


% ────────────────────────────────────────────────────────────
\chapter{Text Feature Construction}
\label{ch:text_features}

\section{Dictionary-Based Features}
\label{sec:dictionary_features}
% DICTIONARIES USED (loaded from text files, all tokens case-folded to lowercase):
%   LM categories: Fin-Neg.txt (negative), Fin-Pos.txt (positive), Fin-Unc.txt (uncertainty),
%     Fin-Lit.txt (litigious), MW-Strong.txt (modal strong), MW-Weak.txt (modal weak).
%   Harvard: Harvard_IV_NEG_Inf.txt (H4N informal negative -- the misclassification baseline).
%   Trading strategy only: fin_neg (financial negative subset).
% TOKENIZATION: Regex [A-Za-z]+(?:'[A-Za-z]+)? -- alphabetic tokens with optional apostrophe
%   contractions (e.g., "don't"). All tokens case-folded. Numbers, punctuation excluded.
% TEXT SCOPES:
%   (1) full_10k: Entire filing text after pre-processing.
%   (2) mda (Item 7): MD&A section only, filtered to item_id="7". Requires token_count >= 250.
% PROPORTIONAL FEATURES: signal_prop = matched_count / document_length.
%   Result in [0, 1]; null if document_length = 0.
% TF-IDF FEATURES (LM2011 Eq. 1 -- implemented exactly):
%   For each token t in signal vocabulary with TF > 0 in document:
%     term_weight(t) = [(1 + log(TF)) / (1 + log(DL))] * log(N / DF)
%   where TF = term frequency, DL = document length (total tokens),
%     N = number of documents in corpus, DF = document frequency of token t.
%   signal_tfidf = sum of term_weight(t) over all signal tokens in document.
%   No sklearn TfidfVectorizer, no smoothing, no +1 offset -- exact LM2011 formula.
% OUTPUT COLUMNS (always _prop and _tfidf pairs):
%   lm_negative_*, lm_positive_*, lm_uncertainty_*, lm_litigious_*,
%   lm_modal_strong_*, lm_modal_weak_*, h4n_inf_*.
%   Trading strategy adds: fin_neg_prop, fin_neg_tfidf.

\section{FinBERT Sentence-Level Sentiment Scoring}
\label{sec:finbert_scoring}
% STATUS: Planned (extension spec E3, status: "planned"). Not yet implemented in code.
% FROM CODE (extension spec): Output schema requires columns:
%   doc_id, filing_date, text_scope, model_name, model_version, segment_policy,
%   score_label, score_value.
%   Grain: (doc_id, text_scope, model_name).
% FROM CODE (extension spec): Design questions still pending:
%   - Scoring grain: sentence, paragraph, item, or filing chunk.
%   - Chunk aggregation from model outputs to doc_id-level features.
%   - Tokenizer, truncation, batching, pooling rules.
%   - Frozen vs. fine-tuned.
% TODO: Describe FinBERT inference setup once implementation decisions are made.

\section{Aggregation from Sentence to Document Level}
\label{sec:aggregation}
% FROM CODE (extension spec): Model-based features must align on doc_id, filing_date,
%   and text_scope with dictionary features. Output includes model_name, model_version,
%   segment_policy columns for traceability.
% FROM CODE (extension spec): Grain is (doc_id, text_scope, model_name).
% TODO: Define aggregation method once FinBERT implementation decisions are made.

\section{Feature Comparison and Descriptive Statistics}
\label{sec:feature_descriptives}
% TODO: Generate from actual data once both dictionary and FinBERT features are computed.
% FROM CODE (extension spec): reporting_splits require replication-overlap window (1994-2008)
%   kept distinct from extension window (2009-2024).
% FROM CODE: Dictionary output columns available for comparison:
%   lm_negative_prop/tfidf, lm_positive_prop/tfidf, lm_uncertainty_prop/tfidf,
%   lm_litigious_prop/tfidf, lm_modal_strong_prop/tfidf, lm_modal_weak_prop/tfidf, h4n_inf_prop/tfidf.


% ────────────────────────────────────────────────────────────
\chapter{Empirical Methodology}
\label{ch:methodology}

\section{Filing-Period Excess Returns}
\label{sec:excess_returns}
% DEPENDENT VARIABLE for return regressions (Tables IV, V, VI).
% EVENT WINDOW: 4 trading days [0, 3] starting from filing_trade_date.
%   filing_trade_date = first trading day on or after the SEC filing calendar date.
% COMPUTATION:
%   event_stock_gross = product(1 + stock_return_t) for t in [0, 3]
%   event_market_gross = product(1 + market_return_t) for t in [0, 3]
%   filing_period_excess_return = event_stock_gross - event_market_gross
% MARKET RETURN: Equal-weighted market from Fama-French daily factors (mkt_rf + rf).
% STOCK RETURN: CRSP daily return including dividends (FINAL_RET, which incorporates
%   delisting returns -- raw RET does not).
% NOTE: This is a buy-and-hold excess return, not a cumulative abnormal return from
%   a market model. Matches LM2011 methodology.

\section{Standardized Unexpected Earnings}
\label{sec:sue}
% DEPENDENT VARIABLE for SUE regressions (Table VIII).
% FORMULA: SUE = (actual_eps - forecast_consensus_mean) / pre_filing_price
%   where actual_eps = reported EPS from Refinitiv/IBES,
%   forecast_consensus_mean = mean analyst EPS forecast,
%   pre_filing_price = stock price on last trading day before filing date.
% ADDITIONAL ANALYST CONTROLS:
%   analyst_dispersion = forecast_std_dev / pre_filing_price
%   analyst_revisions = forecast_revision_4m / prior_month_price
% DATA SOURCE: Refinitiv LSEG API (TR.EPSMean, TR.EPSActValue, TR.EPSStdDev).
%   Matched to filing documents via bridge pipeline (PERMNO+GVKEY -> RIC -> analyst data).

\section{Fama-MacBeth Cross-Sectional Regressions}
\label{sec:fama_macbeth}
% PROCEDURE (Fama-MacBeth 1973):
%   Step 1: For each calendar quarter, run cross-sectional OLS of dependent variable
%     on text signal + control variables + FF48 industry dummies.
%   Step 2: Collect time series of quarterly coefficient estimates.
%   Step 3: Report weighted mean of coefficients (weights = quarterly observation counts).
%   Step 4: Standard errors via Newey-West with 1 lag on the coefficient time series.
%     SE = sqrt(variance) where variance uses Bartlett kernel up to min(nw_lags, T-1).
% CONTROL VARIABLES (return regressions):
%   log_size = log(market_equity) where market_equity = |PRC| * SHROUT on pre_filing_trade_date.
%   log_book_to_market = log(book_equity / market_equity). Book equity uses FF2001 formula:
%     SEQ (or CEQ+PSTK, or AT-LT) adjusted for preferred stock and deferred taxes.
%   pre_ffalpha = intercept from FF3 regression on daily returns over [-252, -6] trading days
%     relative to filing_trade_date. Minimum 60 observations required.
%   log_share_turnover = log(sum of daily volume over [-252, -6] / shares outstanding).
%     Minimum 60 trading days with non-null volume required.
%   nasdaq_dummy = 1 if EXCHCD = 3, else 0.
%   institutional_ownership = percentage from most recent quarter-end before filing date.
% ADDITIONAL CONTROLS (SUE regressions only):
%   analyst_dispersion, analyst_revisions (defined in Section~\ref{sec:sue}).
% INDUSTRY DUMMIES: FF48 classification. First industry omitted (dummy variable trap).
%   Industry codes from SIC via Fama-French 48-industry mapping.
% MINIMUM REQUIREMENTS: pre_filing_price >= \$3.00, 60+ days of pre-filing return data.

\section{Trading Strategy}
\label{sec:trading_strategy}
% REPLICATION WINDOW: July 1, 1997 to June 30, 2007 (strategy returns).
%   Sorting years: 1997-2006 (filings used for quintile assignment).
% SORTING SIGNALS: fin_neg_prop, fin_neg_tfidf, h4n_inf_prop, h4n_inf_tfidf.
% QUINTILE CONSTRUCTION: Within each (sort_year, signal_name), sort filings by signal value.
%   Assign quintiles 1-5 via: quintile = int(rank * 5 / N) + 1.
%   Q1 = highest signal value (most negative/uncertain), Q5 = lowest.
% SORT YEAR: sort_year = filing_date.year + 1. Portfolios formed in July of sort_year.
% HOLDING PERIOD: 12 months (July through June, annual rebalancing).
% PORTFOLIO WEIGHTING: Equal-weighted (default). Value-weighted variant available.
% LONG-SHORT: long_short_return = Q1_return - Q5_return (monthly frequency).
% FF4 ALPHA: Regress monthly long-short returns on four factors:
%   long_short_return = alpha + beta_mkt*MKT_RF + beta_smb*SMB + beta_hml*HML + beta_mom*MOM + epsilon
%   Alpha (intercept) is the risk-adjusted abnormal return of the long-short portfolio.
%   Factor data: Fama-French monthly factors including momentum (from Ken French data library).
%   Factor scaling: if max |factor| > 1, divide by 100 (handles percentage vs. decimal format).


% ────────────────────────────────────────────────────────────
\chapter{Results}
\label{ch:results}

\section{LM2011 Replication (1994--2008)}
\label{sec:replication_results}

\subsection{Return Regressions}
\label{sec:return_regressions}
% FROM CODE (lm2011_regressions.py -- build_lm2011_table_iv_results):
%   TABLE IV (text_scope: full_10k, dep. var: filing_period_excess_return):
%     Four specifications, each one text signal + all controls + FF48 dummies:
%     (1) h4n_inf_prop, (2) lm_negative_prop, (3) h4n_inf_tfidf, (4) lm_negative_tfidf.
% FROM CODE (build_lm2011_table_v_results):
%   TABLE V (text_scope: mda_item_7, dep. var: filing_period_excess_return):
%     Same four specifications as Table IV. Filter: token_count_mda >= 250.
% FROM CODE (build_lm2011_table_vi_results):
%   TABLE VI (text_scope: full_10k, dep. var: filing_period_excess_return):
%     Separate specifications for prop and tfidf variants of all six LM categories:
%     lm_negative, lm_positive, lm_uncertainty, lm_litigious, lm_modal_strong, lm_modal_weak.
% FROM CODE (build_lm2011_table_ia_i_results):
%   Internet Appendix Table IA-I: normalized_difference variants (industry-adjusted).

\subsection{SUE Regressions}
\label{sec:sue_regressions}
% FROM CODE (build_lm2011_table_viii_results):
%   TABLE VIII (text_scope: full_10k, dep. var: sue).
%   Signal columns: h4n_inf_prop, lm_negative_prop, h4n_inf_tfidf, lm_negative_tfidf.
%   Additional controls beyond return set: analyst_dispersion, analyst_revisions.
%   Requires Refinitiv analyst data match (analyst_match_status == "MATCHED").

\subsection{Trading Strategy Returns}
\label{sec:trading_results}
% FROM CODE (build_lm2011_trading_strategy_monthly_returns, build_lm2011_trading_strategy_ff4_summary):
%   Monthly long-short portfolio returns (Q1 - Q5) for each sorting signal.
%   Signals: fin_neg_prop, fin_neg_tfidf, h4n_inf_prop, h4n_inf_tfidf.
%   Period: July 1, 1997 to June 30, 2007 (_STRATEGY_START_DATE, _STRATEGY_END_DATE).
%   FF4 alpha regression on monthly long-short returns (MKT_RF, SMB, HML, MOM).
% FROM CODE (build_lm2011_table_ia_ii_results):
%   Internet Appendix Table IA-II with detailed alpha/beta estimates.

\section{Extended Sample (2009--2024)}
\label{sec:extension_results}
% FROM CODE (extension spec): Run identical specifications on extension_only_window [2009-01-01, 2024-12-31].
%   reporting_splits: replication-overlap window vs. post-2008 extension window kept distinct.
%   "If structural breaks appear material, add narrower subperiod splits instead of pooling
%    silently across the full 1994-2024 horizon." (extension spec)
% FROM CODE (item_regime_10k.json): Regime changes in this window:
%   Item 4 repurposed 2012, Item 6 reserved 2021, Item 9C added 2022, Item 1C added 2023.
% TODO: Discuss results once regressions are run on extended sample.

\section{FinBERT vs.\ Dictionary Feature Comparison}
\label{sec:finbert_results}
% FROM CODE (extension spec): comparison_contract:
%   "Model-based features must be emitted alongside dictionary features on the same sample
%    and scope definitions."
%   "Evaluation should emphasize incremental explanatory power and robustness, not only
%    raw score correlation."
% FROM CODE (extension spec): extension_comparison_results grain:
%   (run_id, sample_window, text_scope, model_name, outcome_name).
%   Required columns: specification_name, estimate, standard_error, p_value, n_obs.
% TODO: Define specific regression specifications once FinBERT features are available.


% ────────────────────────────────────────────────────────────
\chapter{Discussion}
\label{ch:discussion}

\section{Replication Fidelity and Deviations}
\label{sec:replication_fidelity}
% FROM CODE (replication spec): Known scope differences vs. LM2011 paper:
%   - excluded_outputs: Rule 10b-5 fraud label regressions, material-weakness label regressions.
%   - Analyst data sourced via LSEG/Refinitiv API, not original I/B/E/S tape.
%   - FF2001 accounting formulas reimplemented (lm2011.py: book equity hierarchy
%     SEQ > CEQ+PSTK > AT-LT, preferred stock hierarchy PSTKRV > PSTKL > PSTK).
% TODO: Compare actual regression outputs against LM2011 paper tables.

\section{Temporal Stability of Textual Signals}
\label{sec:temporal_stability}
% FROM CODE (extension spec): Results reported separately for 1994-2008 overlap and 2009-2024 extension.
% FROM CODE (item_regime_10k.json): Form structure changes in extension window that may
%   affect token distributions: Item 6 reserved (2021-02-10), Item 1C added (2023-12-15).
% TODO: Discuss temporal stability of coefficients once results are available.

\section{Dictionary vs.\ FinBERT: Complementarity or Substitution?}
\label{sec:dict_vs_finbert}
% FROM CODE (extension spec): Comparison contract positions dictionary as "benchmark
%   representation" and FinBERT as "comparator". Both must share doc_id/text_scope grain.
% TODO: Discuss complementarity vs. substitution once comparison results are available.

\section{Limitations}
\label{sec:limitations}
% FROM CODE: Pipeline processes 10-K and 10-K405 only (lm2011.py form constants);
%   10-Q, 8-K, proxy statements are out of scope.
% FROM CODE: Analyst data sourced via Refinitiv API, not original I/B/E/S tape --
%   coverage may differ. SUE regressions require analyst_match_status == "MATCHED",
%   so SUE sample is a subset of the return regression sample.
% FROM CODE: Extraction uses two-pass strategy with fallback, but non-standard filing
%   formats can still produce zero items (fail-fast with logged context).
% TODO: Discuss further limitations once data review is complete.


% ────────────────────────────────────────────────────────────
\chapter{Conclusion}
\label{ch:conclusion}

\section{Summary of Findings}
\label{sec:summary}
% TODO: Summarize results once regressions are run.
% FROM CODE: Three research axes to address (extension spec):
%   (E1) Horizon Extension, (E2) Additional Item Scopes, (E3) Dictionary vs. FinBERT.

\section{Implications and Future Research}
\label{sec:future_research}
% FROM CODE (extension spec): Candidate item scopes still pending (status: "placeholder"):
%   item_1a_risk_factors, item_7a_market_risk, additional_item_scope_tbd.
%   Prerequisite: "Validate regime-aware availability and extraction quality before activation."
% FROM CODE (extension spec): FinBERT design question still open:
%   "Decide whether to keep the model frozen or allow a later fine-tuning stage."
% TODO: Write implications and future research once results are available.


% ============================================================
% BACK MATTER
% ============================================================
\bibliography{references}

\appendix
\chapter{Regression Variable Definitions}
\label{app:variables}
% TABLE: Complete variable definitions for all regression specifications.
% Columns: Variable Name | Definition | Source | Transformation
%
% filing_period_excess_return | 4-day buy-and-hold excess return [0,3] | CRSP + FF daily | product(1+ret) - product(1+mkt)
% sue | Standardized unexpected earnings | Refinitiv IBES | (actual - forecast) / price
% log_size | Log market equity | CRSP | log(|PRC| * SHROUT) on pre_filing_trade_date
% log_book_to_market | Log book-to-market | Compustat + CRSP | log(BE/ME), FF2001 book equity
% pre_ffalpha | Pre-filing FF3 alpha | CRSP + FF daily | OLS intercept on [-252, -6] trading days
% log_share_turnover | Log share turnover | CRSP | log(sum VOL[-252,-6] / SHROUT)
% nasdaq_dummy | NASDAQ listing indicator | CRSP | 1 if EXCHCD=3, else 0
% institutional_ownership | Institutional holdings % | Refinitiv | Most recent quarter-end before filing
% analyst_dispersion | Forecast std / price | Refinitiv IBES | TR.EPSStdDev / pre_filing_price
% analyst_revisions | 4-month revision / price | Refinitiv IBES | revision_4m / prior_month_price
% lm_negative_prop | LM negative word proportion | SEC text | matched_count / token_count
% lm_negative_tfidf | LM negative TF-IDF | SEC text | LM2011 Eq. (1) weighted sum
% h4n_inf_prop | H4N negative proportion | SEC text | matched_count / token_count

\chapter{Item Regime Change Log}
\label{app:regime}
% TABLE: Complete regime change log from item_regime_10k.json.
% Columns: SEC Release | Year | Item | Change | Effective Date | Trigger Type
%
% 33-7386  | 1997 | 7A  | Added (Market Risk)          | 1997-04-11 | effective_date
% 33-8124  | 2002 | 15  | Controls & Procedures        | 2002-08-29 | effective_date
% 33-8183  | 2003 | 14  | Principal Accountant Fees     | 2003-01-22 | effective_date
% 33-8238  | 2003 | 9A  | Added (Controls/Procedures)   | 2003-08-14 | effective_date
% 33-8400  | 2004 | 9B  | Added (Other Information)     | 2004-08-23 | effective_date
% 33-8591  | 2005 | 1A,1B | Added (Risk Factors, Unresolved Comments) | 2005-12-01 | effective_date
% 33-9089  | 2009 | 4   | Removed (Voting Results)      | 2010-02-28 | effective_date
% 33-9286  | 2011 | 4   | Repurposed (Mine Safety)      | 2012-01-27 | effective_date
% 34-77969 | 2016 | 16  | Added (Form 10-K Summary)     | 2016-11-02 | effective_date
% 33-10890 | 2020 | 6   | Reserved (Selected Fin Data)  | 2021-02-10 | effective_date
% 34-93701 | 2021 | 9C  | Added (HFCAA Disclosure)      | 2022-01-10 | effective_date
% 33-11216 | 2023 | 1C  | Added (Cybersecurity)         | 2023-12-15 | fiscal_year_end_ge

\chapter{FinBERT Inference Configuration}
\label{app:finbert_config}
% To be completed once FinBERT implementation is finalized (extension spec E3).
% Expected contents:
% - Model name, version, and checkpoint hash.
% - Tokenizer: WordPiece, max sequence length 512.
% - Sentence splitting method and library.
% - Batch size, device (GPU model), precision (fp32/fp16).
% - Aggregation method: mean pooling (baseline).
% - Reproducibility: random seed, deterministic mode setting.
% - Runtime: approximate inference time per filing and total corpus.

% --- Statutory Declaration ---
\chapter*{Statutory Declaration}
\addcontentsline{toc}{chapter}{Statutory Declaration}
I hereby declare that I have written this thesis independently and have not used any sources or aids other than those indicated. All passages taken verbatim or in substance from published or unpublished works of others are marked as such. This thesis has not been submitted to any other examination authority in the same or a similar form.

\vspace{2cm}
\noindent Hannover, \today \hfill Signature

\end{document}
